\documentclass[10.5pt,a4paper]{article}
\usepackage[T1]{fontenc}
\usepackage[utf8]{inputenc}
\usepackage{mathptmx} % Times-like serif (Elsevier body font)
\usepackage[margin=1.55cm]{geometry}
\usepackage{xcolor}
\usepackage{enumitem}
\usepackage{titlesec}
\usepackage[protrusion=true,expansion=false]{microtype}
\usepackage[official]{eurosym}
\usepackage[hidelinks]{hyperref}

\definecolor{accent}{HTML}{000000}
\definecolor{accent2}{HTML}{000000}
\definecolor{grey}{HTML}{000000}
\hypersetup{colorlinks=true, urlcolor=accent2, linkcolor=accent2}
\emergencystretch=3em
\setlength{\parindent}{0pt}

\titleformat{\section}{\normalsize\bfseries\color{accent}}{}{0em}{}[{\vspace{-0.3ex}\color{accent}\titlerule[0.8pt]}]
\titlespacing*{\section}{0pt}{1.1ex}{0.5ex}
\setlist[itemize]{leftmargin=1.25em, itemsep=0.5pt, topsep=1pt, parsep=0pt, label=\textbullet}

\newcommand{\dateb}[1]{{\small\itshape\color{grey}#1}}
\newcommand{\cvrow}[2]{%
  \noindent\makebox[2.8cm][l]{\dateb{#1}}%
  \parbox[t]{\dimexpr\linewidth-2.8cm\relax}{#2}\par\smallskip}
\pagestyle{empty}

\begin{document}

{\LARGE\bfseries\color{accent} Dr. NKODIA Hardy Medry Dieu-Veil}\\[2pt]
{\itshape\color{accent2} Geologist — Tectonics, geodynamics and structural geology}\\[3pt]
{\small\color{grey} Brazzaville, Republic of Congo \ \textbullet\ \href{mailto:nkodiahardy@gmail.com}{nkodiahardy@gmail.com}}\\[1pt]
{\small \href{https://orcid.org/0000-0001-6919-6863}{ORCID 0000-0001-6919-6863} \ \textbullet\ \href{https://scholar.google.com/citations?user=gYOI-FkAAAAJ&hl=fr}{Google Scholar} \ \textbullet\ \href{https://www.researchgate.net/profile/Hardy-Medry-Dieu-Veil-Nkodia}{ResearchGate} \ \textbullet\ \href{https://nkodiahardy.com}{nkodiahardy.com}}

\section{Profile}
Lecturer--researcher in structural geology and tectonics: palaeostress analysis, fault reactivation and intraplate seismic hazard in Central Africa. \textit{Interests}: fault/fracture reservoirs (Inkisi, Pool); Congo Basin geomorphology; mineralisations and aquifers; seismic hazard of the Central African rifted margin.

\section{Education}
\cvrow{2019 – 2024}{\textbf{PhD in Geosciences} — Marien Ngouabi University (Congo) \& Royal Museum for Central Africa (Belgium). Highest distinction with the jury's congratulations.}
\cvrow{2015 – 2017}{\textbf{MSc in Geosciences} — Marien Ngouabi University. Distinction; top of class.}
\cvrow{2012 – 2015}{\textbf{BSc in Geosciences} — Marien Ngouabi University. 2nd of class.}
\cvrow{2011}{\textbf{Baccalauréat (series D)} — Lycée Chaminade, Brazzaville.}

\section{Academic appointments}
\cvrow{2025 – present}{\textbf{Assistant} (tenure-track, toward Assistant Professor) — Marien Ngouabi University, FST, Dept. of Geology (Congo). Teaching and research in structural geology; MSc supervision.}
\cvrow{2025 – 2026}{\textbf{Postdoctoral researcher} — Pukyong National University, Busan (South Korea), on research leave from Marien Ngouabi University. Contemporary stress field and seismogenic structures of the Korean Peninsula.}
\cvrow{Aug – Sep 2026}{\textbf{Research visit} — RMCA, Tervuren (Belgium). West Congo Belt and karst.}
\cvrow{Feb–Mar 2025}{\textbf{Postdoctoral researcher} (visit) — Royal Museum for Central Africa (RMCA), Tervuren (Belgium).}
\cvrow{2018 – 2024}{\textbf{Lecturer--Researcher} — Marien Ngouabi University (Congo). Structural Geology course (L2); co-organised an international conference.}
\cvrow{2019 – 2021}{\textbf{Research assistant} — RMCA, Tervuren (Belgium). Managed a \euro{}15,000 budget over 3 years; international collaborations.}
\cvrow{2019 – 2021}{\textbf{Instructor} — African School of Development (EAD), Brazzaville. Structural geology course.}

\section{Grants \& honours}
\begin{itemize}
\item \textbf{AMORCE-GEO (2025–2027)} — project funded by ARES (Belgium); role: co-organiser and trainer.
\item PhD awarded with the jury's congratulations (2024); top of MSc class (2017).
\end{itemize}

\section{Publications — selected (peer-reviewed)}
{\small\color{grey}Full list: 12 peer-reviewed articles (see ORCID / Google Scholar).}
\begin{enumerate}[leftmargin=1.5em, itemsep=2pt, topsep=2pt, label=\arabic*.]
\item \textbf{Nkodia, H.\,M.\,D.-V.}, Hategekimana, F., Naik, S.\,P., Cheon, Y., Peace, A.\,L., Bae, S.-Y., Kandregula, R.\,S., \& Kim, Y.-S. (2026). Vertical kinematic partitioning in intraplate fault systems: evidence from the SE Korean Peninsula. \emph{Journal of Structural Geology}, \textbf{211}, 105778. \href{https://doi.org/10.1016/j.jsg.2026.105778}{doi:10.1016/j.jsg.2026.105778}
\item \textbf{Nkodia, H.\,M.\,D.-V.}, Park, K., Naik, S.\,P., Peace, A.\,L., \& Kim, Y.-S. (2026). Assessing the contemporary stress field and seismogenic structures of the Korean peninsula using focal-mechanism inversion and slip-tendency analysis. \emph{Tectonophysics}, \textbf{934}, 231261. \href{https://doi.org/10.1016/j.tecto.2026.231261}{doi:10.1016/j.tecto.2026.231261}
\item Colet, M., Kolawole, F., Ajala, R., Delvaux, D., \& \textbf{Nkodia, H.\,M.\,D.-V.} (2025). Active crustal deformation across a nucleating extensional microplate, D.\,R. Congo, East Africa. \emph{Tectonics}, 44, e2025TC008815. \href{https://doi.org/10.1029/2025TC008815}{doi:10.1029/2025TC008815}
\item \textbf{Nkodia, H.\,M.\,D.-V.}, Boudzoumou, F., Miyouna, T., \emph{et al.} (2024). Brittle faulting and tectonic stress history on the western margin of the Congo Basin between Kinshasa and Brazzaville. \emph{Tectonophysics}, 877, 230282. \href{https://doi.org/10.1016/j.tecto.2024.230282}{doi:10.1016/j.tecto.2024.230282}
\item Mavoungou, Y.\,B., \textbf{Nkodia, H.\,M.\,D.-V.}, Watha-Ndoudy, N., \& Bolarinwa, A.\,T. (2024). Lineament extraction and paleostress analysis in the Bikélélé iron deposit (Chaillu Massif, Congo). \emph{Modeling Earth Systems and Environment}, 10(2), 1993–2009. \href{https://doi.org/10.1007/s40808-023-01883-3}{doi:10.1007/s40808-023-01883-3}
\item \textbf{Nkodia, H.\,M.\,D.-V.}, Miyouna, T., Kolawole, F., \emph{et al.} (2022). Seismogenic fault reactivation in western Central Africa: insights from regional stress analysis. \emph{Geochemistry, Geophysics, Geosystems}, 23(11), e2022GC010377. \href{https://doi.org/10.1029/2022GC010377}{doi:10.1029/2022GC010377}
\item \textbf{Nkodia, H.\,M.\,D.-V.}, Miyouna, T., Delvaux, D., \& Boudzoumou, F. (2020). Flower structures in sandstones of the Palaeozoic Inkisi Group (Brazzaville, Congo). \emph{South African Journal of Geology}, 123(4), 531–550. \href{https://doi.org/10.25131/sajg.123.0038}{doi:10.25131/sajg.123.0038}
\item Miyouna, T., \textbf{Nkodia, H.\,M.\,D.-V.}, Essouli, O.\,F., \emph{et al.} (2018). Strike-slip deformation in the Inkisi Formation, Brazzaville, Republic of Congo. \emph{Cogent Geoscience}, 4(1), 1542762. \href{https://doi.org/10.1080/23312041.2018.1542762}{doi:10.1080/23312041.2018.1542762}
\end{enumerate}

\section{Invited talk}
\begin{itemize}
\item \textbf{6 July 2026} — ``Ocean--continent tectonic stresses and earthquake hazards of the Central African rifted margin''. Online seminar \emph{Earthquake Hazards of Rifted Margins}, Columbia University / Lamont-Doherty Earth Observatory. \href{https://www.youtube.com/watch?v=lnRf0Jz0QbM}{Recording}.
\end{itemize}

\section{Skills, languages \& service}
\begin{itemize}
\item Tectonic/structural analysis and palaeostress (slip tendency); remote sensing \& GIS (ALOS-PALSAR DEM, satellite/drone imagery); petrology and field geology; Python (publication figures), Adobe Illustrator.
\item \textit{Certifications}: Project Management (Google, 2024); scientific writing/publishing (École Polytechnique Paris, 2023); radar remote sensing (IGNFI/AFD, 2019).
\item \textit{Languages}: French (native), English (professional). \quad \textit{Service}: Founder \& President of Kongo Science; Founder of the LGRN (Marien Ngouabi University).
\end{itemize}

\section{References}
\begin{itemize}
\item \textbf{Prof. F. Boudzoumou} (MNU).
\item \textbf{Dr. D. Delvaux} — Editor-in-Chief, \emph{J. African Earth Sci.}; RMCA, Tervuren.
\item \textit{Referees' contact details available on request.}
\end{itemize}

\end{document}
